% !TEX root = ../QuantumComputing.tex
%% Lecture 3: Quantum Simulation
%% Based on original notes by T.J. Osborne, enriched and unified

\section{Trotterization and quantum simulation}\label{sec:simulation}

\begin{historical}[From Feynman's vision to practical simulators]
  In 1982, Richard Feynman~\cite{Feynman1982} asked whether quantum systems
  could be efficiently simulated on classical computers, and conjectured that
  the answer was no---but that a \emph{quantum} computer could do the job.
  Seth Lloyd~\cite{Lloyd1996} proved Feynman right in 1996, showing that
  local Hamiltonians can be efficiently simulated on a quantum computer via
  the Lie-Trotter product formula.  Quantum simulation remains one of the
  most promising near-term applications of quantum computers, with potential
  impact on quantum chemistry, materials science, and high-energy physics.
\end{historical}

\subsection{Why simulate quantum systems?}\label{subsec:why-simulate}

Complex quantum systems have numerous industrially significant applications,
from molecular structure for quantum chemistry through to materials design.
Despite their importance, progress on understanding the physics of such
systems faces continuing challenges.

The direct simulation of quantum systems by classical computers is possible
in principle but generically \emph{inefficient}.  Consider the Schrödinger
equation for a many-body quantum system:
\begin{equation}\label{eq:schrodinger}
  \eqbox{i\frac{d}{dt}\ket{\psi} = H \ket{\psi}.}
\end{equation}
If one writes this equation in terms of a basis for Hilbert space, one
obtains a system of coupled linear ODEs.  Such systems are straightforward
to solve numerically.  So what is the difficulty?

\begin{intuition}[The exponential wall]
  The key challenge is the \emph{exponential} number of differential
  equations that must be solved.  For one qubit, $2$ equations; for two
  qubits, $4$; and for $n$ qubits, $2^n$ equations.  For a system of just
  $300$ qubits, there are more equations than atoms in the observable
  universe ($2^{300} \approx 10^{90}$).  Sometimes insightful approximations
  reduce the effective number of equations (e.g., tensor network methods for
  low-entanglement states), but for many physically interesting systems no
  such shortcuts are known.  This is one of the strongest motivations for
  building quantum computers.
\end{intuition}

\subsection{Many-body quantum systems}\label{subsec:many-body}

A many-body quantum system is comprised of many quantum degrees of freedom.
We model our system as a collection of $n$ qubits (assuming distinguishable
particles for simplicity) with Hilbert space
\begin{equation}
  \hilb \equiv (\C^2)^{\tp n}.
\end{equation}

\begin{definition}[Local Hamiltonian]\label{def:local-hamiltonian}
  A Hamiltonian $H$ acting on $n$ qubits is called \emph{$c$-local} if it
  can be written as
  \begin{equation}
    H = \sum_{j=1}^L h_j,
  \end{equation}
  where each operator $h_j$ acts nontrivially on at most $c$ qubits (for
  some constant $c$ independent of $n$).  The number $L$ of terms can be
  any polynomial in $n$.
\end{definition}

\begin{example}[Transverse-field Ising model]
  The transverse-field Ising model on a chain of $n$ qubits is
  \begin{equation}
    H = \sum_{j=1}^{n-1} \paulix_j\paulix_{j+1}
    + \lambda \sum_{j=1}^n \pauliz_j,
  \end{equation}
  where $\lambda\in\R$ is the transverse field strength.  This is a
  $2$-local Hamiltonian with $L = 2n-1$ terms.  It models magnetic phase
  transitions and is one of the simplest exactly-solvable interacting
  quantum systems.
\end{example}

\begin{remark}
  The restriction to constant locality $c$ implies that $L$ is bounded
  by a polynomial in $n$.  Indeed, there are at most $\binom{n}{c} = \bigO{n^c}$
  subsets of size $c$, so $L \le \binom{n}{c}$, which is polynomial for
  fixed~$c$.
\end{remark}


\subsection{The quantum simulation algorithm}\label{subsec:sim-algorithm}

The goal of quantum simulation is, for a given time $t$, to prepare the
state
\begin{equation}
  \ket{\psi(t)} = e^{-iHt}\ket{\psi(0)},
\end{equation}
corresponding to the solution of the Schrödinger equation~\eqref{eq:schrodinger}.
The objective is to implement the unitary $U(t) \equiv e^{-iHt}$ with quantum gates.

If we na\"ively applied the universality theorem from Section~\ref{sec:circuits},
we would need roughly $4^n$ elementary gates (since a general $n$-qubit
unitary has $\sim 4^n$ real parameters)---and even computing
a single matrix element of $U(t)$ requires solving the Schrödinger equation
in the first place.  We need a fundamentally different strategy.

\begin{intuition}[The key insight]
  Although $e^{-iHt}$ is hard to implement directly, each factor
  $e^{-ih_j t}$ acts on a small subsystem (at most $c$ qubits) and is
  therefore straightforward to compile into a quantum circuit using $\bigO{1}$
  gates.  The challenge is that in general $\comm{h_j}{h_k}\neq 0$, so
  $e^{-iHt} \neq \prod_{j=1}^L e^{-ih_j t}$.  The Lie-Trotter formula
  provides a systematic way to overcome this obstacle.
\end{intuition}

\begin{proposition}\label{prop:commuting-case}
  If $\comm{h_j}{h_k}= 0$ for all $j,k$, then
  \begin{equation}
    U(t) = e^{-iHt} = \prod_{j=1}^L e^{-ih_jt}
  \end{equation}
  for all $t$.
\end{proposition}
\begin{proof}
  It suffices to show $e^{A+B}=e^A e^B$ when $\comm{A}{B}=0$; the general
  case follows by induction on $L$.  Define the matrix-valued functions
  $f(s) = e^{(A+B)s}$ and $g(s) = e^{As}\,e^{Bs}$.  Both satisfy $f(0)=g(0)=\I$.
  Differentiating:
  \[
    f'(s) = (A+B)\,f(s).
  \]
  For $g$, since $\comm{A}{B}=0$ implies $\comm{A}{e^{Bs}}=0$ (expand
  $e^{Bs}$ as a power series; $A$ commutes with every term $B^k$), we get
  \[
    g'(s) = A\,e^{As}\,e^{Bs} + e^{As}\,B\,e^{Bs}
    = (A+B)\,g(s).
  \]
  Since $f$ and $g$ satisfy the same linear ODE $y'=(A+B)y$ with the same
  initial condition $y(0)=\I$, uniqueness of solutions gives $f(s)=g(s)$
  for all $s$.  Setting $s=t$ and $A=-ih_j t$, $B=-i\sum_{k\neq j}h_k t$
  (with appropriate rescaling) yields the result.
\end{proof}

\subsubsection{The Lie-Trotter formula}

\begin{theorem}[Lie-Trotter formula]\label{thm:lie-trotter}
  Let $A$ and $B$ be Hermitian operators.  Then for any real $t$,
  \begin{equation}
    \eqbox{\lim_{n\rightarrow \infty}
    \left(e^{iA\frac{t}{n}}e^{iB\frac{t}{n}}\right)^{n}
    = e^{i(A+B)t}.}
  \end{equation}
\end{theorem}

\begin{proof}
  Write $\delta = t/n$.  By Taylor expansion to second order,
  \begin{equation}
    e^{iA\delta} = \I + iA\delta - \tfrac{1}{2}A^2\delta^2
    + \bigO{\norm{A}^3\delta^3},
  \end{equation}
  and similarly for $e^{iB\delta}$.  Multiplying:
  \begin{equation}\label{eq:trotter-single-step}
    e^{iA\delta}e^{iB\delta} = \I + i(A+B)\delta
    - \tfrac{1}{2}(A^2+2AB+B^2)\delta^2
    + \bigO{\delta^3}.
  \end{equation}
  Meanwhile, $e^{i(A+B)\delta} = \I + i(A+B)\delta
  - \frac{1}{2}(A+B)^2\delta^2 + \bigO{\delta^3}$.
  The difference is
  \begin{equation}\label{eq:trotter-local-error}
    e^{i(A+B)\delta} - e^{iA\delta}e^{iB\delta}
    = \tfrac{1}{2}\comm{A}{B}\,\delta^2 + \bigO{\delta^3},
  \end{equation}
  so
  $\norm{e^{i(A+B)\delta} - e^{iA\delta}e^{iB\delta}}
  \le C\,\delta^2$ for a constant $C$ depending on $\norm{A}$, $\norm{B}$.

  To pass from one step to $n$ steps, use the \emph{telescoping identity}:
  for any operators $P$ and $Q$,
  \begin{equation}\label{eq:telescope}
    P^n - Q^n = \sum_{k=0}^{n-1} P^{n-1-k}(P-Q)Q^k.
  \end{equation}
  Since both $P = e^{i(A+B)\delta}$ and $Q = e^{iA\delta}e^{iB\delta}$ are
  unitary (hence $\norm{P}=\norm{Q}=1$), taking norms gives
  \begin{equation}
    \norm{e^{i(A+B)t} - \left(e^{iA\frac{t}{n}}e^{iB\frac{t}{n}}\right)^n}
    \le n\cdot C\,\delta^2 = \frac{C\,t^2}{n}
    \xrightarrow{n\to\infty} 0.
    \qedhere
  \end{equation}
\end{proof}

By inductively applying the Lie-Trotter formula, one demonstrates that for
$n$ large enough,
\begin{equation}
  U(t) = e^{-it\sum_{j=1}^L h_j}
  \approx \left(e^{-i h_1\frac{t}{n}}e^{-i h_2\frac{t}{n}}
  \cdots e^{-i h_L\frac{t}{n}}\right)^{n}.
\end{equation}
Notice that this is now a product of \emph{local} unitary operators, each
computable with $\bigO{1}$ gates.

\begin{keyresult}[First-order Trotter error bound]
  The first-order Trotter approximation satisfies
  \begin{equation}\label{eq:trotter-error-bound}
    \norm{e^{-iHt} - \left(\prod_{j=1}^L e^{-ih_j t/n}\right)^n}
    \le \frac{t^2}{2n}\sum_{j<k}\norm{\comm{h_j}{h_k}}.
  \end{equation}
  Writing $h_{\max}\equiv \max_j\norm{h_j}$ and bounding
  $\sum_{j<k}\norm{\comm{h_j}{h_k}} \le 2\binom{L}{2}h_{\max}^2
  \le L^2 h_{\max}^2$, accuracy $\epsilon$ requires
  $n = \bigO{t^2 L^2 h_{\max}^2/\epsilon}$ Trotter steps.  Each step
  comprises $L$ local unitaries, giving a total gate count of
  $\bigO{t^2 L^3 h_{\max}^2/\epsilon}$.
\end{keyresult}

\begin{proof}[Proof sketch]
  \textbf{Single-step error.}
  For two operators, the calculation in~\eqref{eq:trotter-local-error}
  gives
  $\norm{e^{-i(A+B)\delta} - e^{-iA\delta}e^{-iB\delta}}
  \le \frac{\delta^2}{2}\norm{\comm{A}{B}} + \bigO{\delta^3}$.
  This extends to $L$ operators by induction: writing
  $H_k = \sum_{j=1}^k h_j$ and peeling off one term at a time,
  \begin{align}
    \norm{e^{-iH\delta} - \prod_{j=1}^L e^{-ih_j\delta}}
    &\le \sum_{k=1}^{L-1}
    \norm{e^{-iH_k\delta}e^{-ih_{k+1}\delta} - e^{-i(H_k+h_{k+1})\delta}}
    \nonumber\\
    &\le \frac{\delta^2}{2}\sum_{k=1}^{L-1}\norm{\comm{H_k}{h_{k+1}}}
    + \bigO{\delta^3}.
  \end{align}
  Since $\norm{\comm{H_k}{h_{k+1}}} \le \sum_{j=1}^k \norm{\comm{h_j}{h_{k+1}}}$,
  the single-step error is at most
  $\frac{\delta^2}{2}\sum_{j<k}\norm{\comm{h_j}{h_k}}$ (up to
  higher-order terms).

  \textbf{Accumulation over $n$ steps.}
  By the telescoping identity~\eqref{eq:telescope} with unitary factors
  ($\norm{\cdot}=1$), the total error after $n$ steps with
  $\delta=t/n$ is at most $n$ times the single-step error:
  \[
    n\cdot \frac{(t/n)^2}{2}\sum_{j<k}\norm{\comm{h_j}{h_k}}
    = \frac{t^2}{2n}\sum_{j<k}\norm{\comm{h_j}{h_k}}.
    \qedhere
  \]
\end{proof}

\subsubsection{The quantum simulation algorithm}

\begin{algorithmbox}[Quantum simulation via Trotterization]
  \begin{itemize}
    \item \textbf{Input}: (1) A Hamiltonian $H = \sum_{j=1}^L h_j$ acting
      on $n$ qubits, where each $h_j$ acts on a constant-size subsystem;
      (2) an initial state $\ket{\psi_0}$; (3) a target accuracy $\delta>0$;
      and (4) a final time $t_f$.
    \item \textbf{Output}: A state $\ket{\widetilde{\psi}(t_f)}$ such that
      $|\braket{\widetilde{\psi}(t_f)}{e^{-iHt_f}\psi_0}|^2 \ge 1-\delta$.
    \item \textbf{Runtime}: $\bigO{\poly{1/\delta}}$ operations.
    \item \textbf{Procedure}: Define the Trotter step
      $U_{\Delta t} \equiv e^{-i h_1\Delta t}e^{-i h_2\Delta t}
      \cdots e^{-i h_L\Delta t}$ and choose $\Delta t$ for acceptable error:
      \begin{enumerate}
        \item Initialise: $\ket{\psi_0}$; set $j=0$.
        \item Update: $\ket{\widetilde{\psi}_{j+1}}
          = U_{\Delta t}\ket{\widetilde{\psi}_{j}}$.
        \item Loop: $j \leftarrow j+1$; repeat step~2 until
          $j\cdot\Delta t \ge t_f$.
        \item Output: $\ket{\widetilde{\psi}(t_f)}
          =\ket{\widetilde{\psi}_{j}}$.
      \end{enumerate}
  \end{itemize}
\end{algorithmbox}

\begin{intuition}[Trotterization as time-slicing]
  Trotterization is analogous to the operator-splitting methods used in
  classical numerical analysis.  By dividing the total evolution time into
  many small steps $\Delta t$, the error from non-commutativity is controlled:
  during each small step, the terms ``almost commute'' because their
  commutator is multiplied by $\Delta t^2$.  The total circuit looks like:

  \begin{center}
  \begin{tikzpicture}[scale=0.8]
    % Time axis
    \draw[->,spacecadet,thick] (0,0) -- (14,0) node[right,font=\small] {$t$};
    % Time slices
    \foreach \i in {0,...,6} {
      \draw[spacecadet] ({2*\i},0.15) -- ({2*\i},-0.15);
    }
    \node[below,font=\scriptsize] at (0,-0.15) {$0$};
    \node[below,font=\scriptsize] at (2,-0.15) {$\Delta t$};
    \node[below,font=\scriptsize] at (4,-0.15) {$2\Delta t$};
    \node[below,font=\scriptsize] at (12,-0.15) {$t_f$};
    \node at (8,-0.15) {$\cdots$};
    % Trotter steps
    \foreach \i in {0,...,2} {
      \node[above,font=\scriptsize,cgblue] at ({2*\i+1},0.2) {$U_{\Delta t}$};
    }
    \foreach \i in {5} {
      \node[above,font=\scriptsize,cgblue] at ({2*\i+1},0.2) {$U_{\Delta t}$};
    }
  \end{tikzpicture}
  \end{center}
\end{intuition}


\subsection{Higher-order product formulas}\label{subsec:higher-order}

The first-order Trotter formula can be improved.  The \emph{second-order
  Suzuki-Trotter} formula uses the \emph{symmetric} (palindromic) product
\begin{equation}
  S_2(\Delta t) = e^{-ih_1 \Delta t/2}e^{-ih_2 \Delta t/2}\cdots
  e^{-ih_L \Delta t}\cdots e^{-ih_2 \Delta t/2}e^{-ih_1 \Delta t/2},
\end{equation}
which achieves error $\bigO{\Delta t^3}$ per step (compared to
$\bigO{\Delta t^2}$ for first order), giving a total error of
$\bigO{t^3/n^2}$ after $n$ steps.

\begin{intuition}[Why symmetrization helps]
  The key observation is that the first-order error is
  $e^{-iA\delta}e^{-iB\delta} = e^{-i(A+B)\delta -\frac{1}{2}\comm{A}{B}\delta^2 + \cdots}$.
  Reversing the order gives
  $e^{-iB\delta}e^{-iA\delta} = e^{-i(A+B)\delta +\frac{1}{2}\comm{A}{B}\delta^2 + \cdots}$.
  The leading error terms have \emph{opposite signs}.  By taking half-steps
  in each order and combining them symmetrically, the $\delta^2$ errors
  cancel exactly, leaving only $\bigO{\delta^3}$ corrections.  This is
  the operator analogue of the midpoint rule improving on the Euler method
  in numerical ODE solving.
\end{intuition}

Even higher-order formulas can be constructed recursively: Suzuki showed
that $S_{2k}(\delta) = [S_{2k-2}(p_k\delta)]^2\, S_{2k-2}((1-4p_k)\delta)\,
[S_{2k-2}(p_k\delta)]^2$ with $p_k = (4-4^{1/(2k-1)})^{-1}$ achieves
error $\bigO{\delta^{2k+1}}$ per step~\cite{ChildsTrotterError2021}.


\subsection{Advanced techniques and further reading}\label{subsec:advanced}

With the deployment of near-term quantum devices, new possibilities for
quantum simulation have emerged.  Beyond Lie-Trotter expansions, the major
modern approaches include:

\begin{itemize}
  \item \textbf{Quantum signal processing and quantum singular value
      transform (QSVT)}~\cite{QSVT2019}: A unifying framework that achieves
    near-optimal Hamiltonian simulation and encompasses many quantum
    algorithms as special cases.

  \item \textbf{Qubitization}~\cite{LowChuang2019}: A quantum walk-based
    method that achieves optimal scaling in all parameters simultaneously.

  \item \textbf{Linear combination of unitaries
      (LCU)}~\cite{ChildsLCU2012}: Expresses $e^{-iHt}$ as a linear
    combination of easily-implementable unitaries, using an ancilla register
    and amplitude amplification.

  \item \textbf{Randomised methods (qDRIFT)}~\cite{Campbell2019}: Randomly
    samples terms from the Hamiltonian at each step, achieving gate
    complexity independent of the number of terms~$L$.
\end{itemize}

\begin{remark}
  The optimal query complexity for Hamiltonian simulation is
  $\Theta(t\norm{H} + \log(1/\epsilon)/\log\log(1/\epsilon))$, achieved by
  QSVT-based methods.  The Trotter approach, while not optimal, remains
  practically competitive due to its simplicity and favourable constant
  factors.
\end{remark}
